% Hand-written sentences that depend on results; every number is a macro from v17_numbers.tex (Research/code/v17_tables.py).
\newcommand{\nPointerSentence}{at $b\in\{3,4\}$ fixed forms lose exactly $4$ and $8$, two rounds lose $1.80$ and between $3.32$ and $3.33$, against the upper bounds $2$ and $4$ of part (a), and $b-1$ rounds lose $1.80$ at both caps (Appendix~\ref{app:pointerexact}).}
\newcommand{\nOrderSentence}{in an illustrative response model, a 10\% chance that the last answer flips biases branching and form scores by a median \nPosBranch{} and \nPosFit{} branching standard errors at $n=1000$ (Appendix~\ref{app:measurement}).}
\newcommand{\nEmAppendixText}{At $n=1000$ the working model on the fitted-form records has the smallest RMSE of the four estimators in most settings, \nEmBrLo--\nEmBrHi{} times that of branching: the model shares strength across paths that the design-based scores treat separately. Its bias is between \nEmBiasLo{} and \nEmBiasHi{} branching standard deviations and its coverage \nEmCovLo--\nEmCovHi{} (below 0.925 in \nEmBelow{} of \nEmSettings{} settings). The bias is a property of the model, not of the sample size: the accuracy of a depth-$b$ tree depends on interactions of up to $b$ items and the label, which a two-way model does not represent. At $n=\nEmBigN$ the same bias is up to \nEmBigBiasHi{} of the estimator\textquotesingle s own standard deviations, coverage falls to \nEmBigCovLo--\nEmBigCovHi{} (below 0.90 in \nEmBigBelow{} of \nEmBigSettings{} settings), and the RMSE is \nEmBigRmseLo--\nEmBigRmseHi{} times that of branching and above that of the design-based form score in \nEmBigWorse{} settings; the design-based scores keep coverage \nEmBigFixCovLo--\nEmBigFixCovHi. A richer model would reduce the bias and raise the variance; with all interactions it is the design-based estimator. On balanced random forms the model-based estimate is unbiased to a similar degree but \nEmUnifLo--\nEmUnifHi{} times less accurate at $n=1000$ than the design-based score on optimized forms, so the choice of forms matters for model-based analysis as well.}
\newcommand{\nDiscussionSentence}{On the trained rules we examined, optimized forms needed between $1.0$ and \nMaxAll{} times the records of branching without a model, one adaptive round recovered at least \nTwoMinPct\% of the log-gap wherever its optimum was computed, designs learned from a smoothed pilot were within \nSmoothOverLo--\nSmoothOverHi\% of the oracle for audits of 8,000 records, and under a pairwise model the price of pre-commitment fell to \nModelTwoLo--\nModelTwoHi.}
\newcommand{\nAbstractData}{On screening trees and forest updates trained on survey and clinical data, optimized forms need $1.0$--\nMaxAll{} times the records of branching without a model, \nSmoothRunLo--\nSmoothRunHi{} when the design is learned from a pilot of the same audit, and \nModelTwoLo--\nModelTwoHi{} under a fitted pairwise model; in simulated audits of up to 8,000 respondents, pairwise maximum likelihood is about as accurate on pre-committed forms as on branching records, while a three-way model and the design-based scores pay \nAudThreeFormsLo--\nAudDesignHi{} in mean squared error.}
\newcommand{\nContribData}{On screening trees and forest updates, optimized forms need $1.0$--\nMaxAll{} times the records of branching without a model and \nModelTwoLo--\nModelTwoHi{} under a pairwise model; a second round recovers most of the design-based gap; and in simulated audits the price of pre-commitment rises with the interactions the analysis allows, from about one for pairwise maximum likelihood to \nAudDesignLo--\nAudDesignHi{} for design-based scores.}
\newcommand{\nAuditParagraph}{Figure~\ref{fig:rounds}(b) and Table~\ref{tab:curve} follow \nCurveTrees{} trees at caps four and five through audits of 500 to 32,000 respondents collected in three ways under the same cap: fitted forms, forms optimized for the pairwise model, and branching, which administers the rule. The price of pre-commitment grows with the interactions that the analysis allows, as Theorem~\ref{thm:structured} predicts. For the design-based scores the mean squared error of forms is \nAudDesignLo--\nAudDesignHi{} times that of branching (medians over trees at each size); for maximum likelihood under the three-way model, forms against branching records give \nAudThreeFormsLo--\nAudThreeFormsHi; under the pairwise model they give \nAudTwoFormsSurveyLo--\nAudTwoFormsSurveyHi{} up to 8,000 respondents, and \nAudTwoMformsSurveyLo--\nAudTwoMformsSurveyHi{} when the forms are optimized for the model. At these sizes pairwise maximum likelihood on forms matches the design-based mean under branching (\nAudTwoVsBranchSurveyLo--\nAudTwoVsBranchSurveyHi). The model is not free, because its bias does not shrink. At 32,000 respondents the pairwise fit is off by a median \nAudBiasFormsBig{} of its standard deviation on forms and \nAudBiasBranchBig{} on branching records; it is less accurate than the design-based form score for \nCurveTwoWorsed{} of the \nCurveTrees{} trees on forms and than the plain mean of $t$ for \nAudBranchTwoWorsed{} on branching records, the ratio between the two designs reaches \nAudTwoFormsMaxAll{} for the worst tree, and its intervals under-cover (Table~\ref{tab:em}). At survey sizes a pairwise model on pre-committed forms thus gives up little to adaptivity; the constants without a model price pre-commitment for claims that must not rest on a model, for richer models and for large audits.}
\newcommand{\nAuditAppendixText}{For the design-based scores the ratio of mean squared errors stays near the population ratio at every size (medians \nAudDesignLo--\nAudDesignHi). For the three-way model the ratio between forms and branching records is \nAudThreeFormsLo--\nAudThreeFormsHi{} on the fitted forms and \nAudThreeMformsLo--\nAudThreeMformsHi{} on the forms optimized for the pairwise model; for the pairwise model it is \nAudTwoFormsSurveyLo--\nAudTwoFormsSurveyHi{} and \nAudTwoMformsSurveyLo--\nAudTwoMformsSurveyHi{} up to 8,000 respondents (at most \nAudTwoFormsSurveyMax{} and \nAudTwoMformsSurveyMax{} for a single tree). These orders of magnitude agree with the population constants of Table~\ref{tab:model}, although the target laws do not satisfy the models. Branching here is the rule's own tree, not the branching design that is optimal within the model, which Table~\ref{tab:model} puts a few percent below the optimal forms. At large sizes the comparison is governed by the bias of the model, on either kind of record: at 32,000 respondents the pairwise fit has a median absolute bias of \nAudBiasFormsBig{} of its standard deviation on forms (\nAudBiasFormsSmall{} at 1,000) and \nAudBiasBranchBig{} on branching records (\nAudBiasBranchSmall{} at 1,000), the pairwise fit to branching records is less accurate than the plain mean of $t$ for \nAudBranchTwoWorsec{} and \nAudBranchTwoWorsed{} of the \nCurveTrees{} trees at 16,000 and 32,000, and the ratio between the two designs ranges from \nAudTwoFormsMinAll{} to \nAudTwoFormsMaxAll{} over trees. Against the design-based form score on the same fitted-form records, pairwise maximum likelihood has the smaller mean squared error for every tree up to $n=\nCurveLastAllBetter$ and the larger for \nCurveTwoWorsec{} and \nCurveTwoWorsed{} of the \nCurveTrees{} trees at 16,000 and 32,000 (RMSE up to \nCurveTwoMaxc{} and \nCurveTwoMaxd{} times larger); which trees these are cannot be read from the capped records. A model-assisted score, which uses a working model only to choose the tables of a pointwise-unbiased score, as in the smoothed pilot of Appendix~\ref{app:pilot}, cannot have a smaller variance than the oracle tables for the same form distribution; the advantage of maximum likelihood at small $n$ is bought with the model's bias.}
\newcommand{\nForestScaleSentence}{On twenty Spambase features the picture is the same at $b=3$: forms need \nFScaleThreeLo--\nFScaleThreeHi{} times the records of branching and two rounds at most \nFScaleThreeTwoHi{} (Table~\ref{tab:forestscale}).}
\newcommand{\nForestScaleAppendixText}{At $b=3$ column generation over branching interviews closed within its budget on all three splits, so the ratios \nFScaleThreeLo--\nFScaleThreeHi{} are exact to the certified gap; they lie in the range of the twelve-feature tasks. The two-round programs did not close in 15 minutes, and their certified intervals put two rounds between \nFScaleThreeTwoLo{} and \nFScaleThreeTwoHi{} times the records of branching. At $b=4$ (4,845 forms) the fixed optimum was certified in seven minutes, but 25 minutes of column generation over branching interviews left only the trivial lower bound for the one split that was run, so its interval is uninformative; the other two splits were not run. Exact adaptive designs at $m=20$ and $b\ge4$ are beyond this implementation.}
